\documentclass[10pt,a4paper]{amsart}
\usepackage{amsmath,amssymb,amsthm,graphicx}
\RequirePackage{amssymb,amsmath,amsthm,bm}
\RequirePackage{thmtools}
\RequirePackage[capitalize]{cleveref}
\providecommand{\tightlist}{%
  \setlength{\itemsep}{0pt}\setlength{\parskip}{0pt}}

%%%%% varphi instead of phi
\renewcommand{\phi}{\varphi}

\mathchardef\mhyph="2D 
\providecommand{\A}{\mathcal{A}}
\providecommand{\C}{\mathcal{C}}
\providecommand{\L}{\mathcal{L}}
\providecommand{\B}{\mathcal{B}}
\providecommand{\M}{\mathcal{M}}
\renewcommand{\S}{\mathcal{S}}
\providecommand{\ol}[1]{\overline{#1}}
\providecommand{\mbf}[1]{\mathbf{#1}}
\providecommand{\mf}[1]{\mathfrak{#1}}
\providecommand{\mrm}[1]{\mathrm{#1}}
\providecommand{\mc}[1]{\mathcal{#1}}
\providecommand{\ra}{\rightarrow}
\providecommand{\Ra}{\Rightarrow}
\providecommand{\La}{\Leftarrow}
\providecommand{\lr}{\leftrightarrow}
\providecommand{\hf}{\mathbb{HF}}
\providecommand{\spec}{DgSp}
\providecommand{\bSigma}{\bm\Sigma}
\providecommand{\bDelta}{\bm\Delta}
\providecommand{\bPi}{\bm\Pi}
\providecommand{\Pinf}[1]{\Pi^{\mathrm{in}}_{#1}}
\providecommand{\Sinf}[1]{\Sigma^{\mathrm{in}}_{#1}}
\providecommand{\dSinf}[1]{d\mhyph\Sigma^{\mathrm{in}}_{#1}}
\providecommand{\Dinf}[1]{\Delta^{\mathrm{in}}_{#1}}
\providecommand{\bp}{\bar p}
\providecommand{\bq}{\bar q}
\providecommand{\ba}{\bar a}
\providecommand{\bb}{\bar b}
\providecommand{\bc}{\bar c}
\providecommand{\bd}{\bar d}
\providecommand{\bx}{\bar x}
\providecommand{\by}{\bar y}
\providecommand{\bz}{\bar z}
\providecommand{\LO}{\mathrm{LO}}
\declaretheorem[name={Theorem},numberwithin=section]{theorem}
\declaretheorem[name={Proposition}, sibling=theorem]{proposition}
\declaretheorem[name={Lemma}, sibling=theorem]{lemma}
\declaretheorem[name={Fact}, sibling=theorem]{fact}
\declaretheorem[name={Claim}, numberwithin=theorem]{claim}
\declaretheorem[name={Definition},numberwithin=section,sibling=theorem, style=definition]{definition}
\declaretheorem[name={Remark}, numbered=no,style=definition]{remark}
\declaretheorem[name={Corollary}, sibling=theorem]{corollary}
\declaretheorem[name={Question},sibling=theorem]{question}
\newcommand{\T}{{\mathcal{T}}}
\newcommand{\SR}{SR}
\newcommand{\pSR}{pSR}
\newcommand{\Q}{\mathbb{Q}}
\newcommand{\Z}{\mathbb{Z}}
\newcommand{\K}{\mathcal{K}}
\newcommand{\N}{\mathcal{N}}
\renewcommand{\L}{\mathcal{L}}
\providecommand{\Ginf}[1]{\Gamma^{\mathrm{in}}_{#1}}
\title{Linear orders of successor Scott sentence complexities exist in all three complexity classes}
\author{David Gonzalez, Dino Rossegger}
\date{arxiv:2305.07126v1; CC BY 4.0}
\begin{document}
\maketitle
\noindent\textit{Editorial note.} The counted statement is the author environment at source0.tex lines 1361-1368 and its proof environment opens at line 1370, the line immediately after; the pairing is by adjacency and the retained proof range 1370-1392 is the whole environment.\par\medskip
\noindent\textit{Editorial note.} The retained preamble is the author style file dino.sty (package lines, the theorem declaration support, the notation and tuple macros and the theorem environments) followed by the ordinal and rank macros of the manuscript itself, reproduced verbatim; this is the same recipe as the previously published card of this article, so no package outside it is requested and extra\_packages stays empty.\par\medskip
\noindent\textit{Editorial note.} Metadata provenance: source\_id and source\_version are the values recorded in this source licence evidence file (sources/scott-linear-orders-2305.07126v1/LICENSE\_EVIDENCE.json: arxiv:2305.07126v1 / 2305.07126v1), which the card points at directly; the title and author list are transcribed from the authors\_index and title keys of the source record of the pinned copy. The card carries no counter setting and imposes no original numbering, since this article has no label-number dump and its numbering cannot be reconstructed from the retained ranges alone.\par\medskip
\noindent\textit{Editorial note.} Bibliography: the retained proof cites ash1986 twice and AR20 once, and the article own bibliography extract staged beside the source (sources/scott-linear-orders-2305.07126v1/limit-plus-three-bibliography.tex) carries both entries verbatim together with ash2000 and montalban2021; it is attached through bibliography\_paths and expanded as the author wrote it. No entry was edited, added or rewritten, and no citation is supplied from knowledge.\par\medskip
\noindent\textit{Editorial scope.} One original result of the pinned author manuscript, counted once: the existence of linear orders in every successor Scott sentence complexity of the three classes (author label cor:SSCsuccessors).\par\medskip
\noindent\textit{Allowed standard prerequisites (editorial).}\begin{itemize}
\item Successor complexity classes. S\_infinity, dS\_infinity and P\_infinity denote the successive classes of Scott sentence complexity, and complexity is the least rank at which a Scott sentence exists.. They are the three families of complexities realised.
\item Multiplication of orders by Z to the alpha. For an order L the order Z\textasciicircum{}alpha times L has the back-and-forth relations obtained from those of L by shifting by twice alpha.. It shifts the known examples of complexity S\_infinity 4 and 5 to every successor rank.
\item Back-and-forth analysis of well-orderings. The back-and-forth relations of well-orderings are known and give the optimality of the Scott sentence complexities of omega to the alpha and of omega to the alpha times two.. It supplies the two base families with odd successor ranks.
\end{itemize}\medskip
\begin{corollary}\label{cor:SSCsuccessors}
There are linear orders of the following Scott sentence complexities:
\begin{enumerate}
\item $\Sinf{\alpha+n}$ for any countable ordinal $\alpha$ and $n\geq 4$,
\item $\dSinf{\alpha+n}$ for any countable ordinal $\alpha $ and $n\geq 1$,
\item $\Pinf{\alpha+n}$ for any countable ordinal $\alpha$ and $n\geq 1$.
\end{enumerate}
\end{corollary}

\begin{proof}
Let $\L_4$ be the constructed example of Scott sentence complexity $\Sinf4$ and $\K_4$ such that $\L_4\leq_4 \K_4$. 
Similarly define $\L_5$ and $\K_5$.
For any $\alpha$ we know that $\Z^\alpha\cdot \L_4$ has a $\Sinf{2\alpha+4}$
Scott sentence, yet it has no $\Pinf{2\alpha+4}$ Scott sentence as
$\Z^\alpha\cdot \L_4\leq_{2\alpha+4} \Z^\alpha\cdot \K_4$. Thus, $\Z^\alpha\cdot
\L_4$ has Scott sentence complexity $\Sinf{2\alpha+4}$. Similarly, $\Z^\alpha\cdot \L_5$ has Scott sentence complexity $\Sinf{2\alpha+5}$. These constructions provide examples for all of the claimed cases of the form $\Sinf{\alpha+n}$.

Let $\L_2:=\Q+2+\Q$. It is not difficult to see that this order has a $\dSinf2$
Scott sentence and that $\Q+3+\Q\leq_2 \L_2\leq_2 \Q$, so this is indeed
optimal. Using the same reasoning as above, we see that $\Z^\alpha\cdot \L_2$
has Scott sentence complexity $\dSinf{2\alpha+2}$. Ash~\cite{ash1986} analyzed
the back-and-forth relations of well-orderings and it follows from his results
that $\omega^\alpha\cdot 2$ has Scott sentence complexity $\dSinf{2\alpha+1}$
(see~\cite[proof of Proposition 19]{AR20}). These constructions provide examples
for all of the claimed cases of the form $\dSinf{\alpha+n}$.

It is a basic exercise to show that $\Q$ has Scott sentence complexity $\Pinf2$.
Using the same reasoning as above, we see that $\Z^\alpha\cdot \Q$ has Scott
sentence complexity $\Pinf{2\alpha+2}$. Furthermore, it again follows from
results of Ash~\cite{ash1986} that $\omega^\alpha$ has Scott sentence complexity
$\Pinf{2\alpha+1}$. These constructions provide examples for all of the claimed
cases of the form $\Pinf{\alpha+n}$. \end{proof}

\begin{thebibliography}{9}
\bibitem{ash2000} C. Ash and J. Knight. \emph{Computable Structures and the Hyperarithmetical Hierarchy}. Studies in Logic and the Foundations of Mathematics144,2000.
\bibitem{AR20} R. Alvir and D. Rossegger. The complexity of Scott sentences of scattered linear orders. \emph{Journal of Symbolic Logic}85(3)(2020),1079--1101.
\bibitem{ash1986} C.J. Ash. Recursive labelling systems and stability of recursive structures in hyperarithmetical degrees. \emph{Transactions of the American Mathematical Society}298(2)(1986),497--514.
\bibitem{montalban2021} A. Montalban. \emph{Computable Structure Theory: Beyond the Arithmetic}. Draft,2021. \texttt{https://math.berkeley.edu/\string~antonio/CSTpart2\_DRAFT.pdf}.
\end{thebibliography}
\end{document}
